\documentclass{article}
\usepackage[utf8]{inputenc}
\usepackage[T1]{fontenc}
\usepackage[margin=2cm]{geometry}
\usepackage{hyperref}
\sloppy
\begin{document}
\section*{Historical source}
Encyclopaedia Britannica scan, source row 104. Parse 3.5 extracted the text. Super 3.5 VL corrected it from an enlarged crop of each block. It can still contain recognition errors.
Page headers, footers, and page numbers are removed.

\subsection*{WEIGHTS AND MEASURES (E0002)}
\textit{Continues an entry from an earlier page.}

litre or cubic decimetre = 10 decilitres = 100 centilitres
27052 Brit. cubic inches; and 50 litres are nearly 11 Brit.
Also 100 litres = 10 decalitres = 1 hectolitre = 2751207
bushels; and 32 hectolitres are nearly equal to 11 Brit.
quarts, and 100 hectolitres = 10 kilolitres or cubic metres = 1
kilolitre = 34390086 Brit. quarters.

The gramme is a weight equal that of 1 cubic centimetre of water at its maximum density. It contains 10 decigrammes = 100 centigrammes = 15.434 troy grains; and 1000 grammes = 100 decigrammes = 10 hectogrammes = 1 kilogramme = 2.679514 lbs. troy = 2.204857 lbs. avoirdupois. Also 100 kilogrammes = 10 myriagrammes = 1 metrical quintal = 220.488 lbs. avoirdupois; and 10 quintals = weight of a cubic metre of water = 1 millier or marine ton = 19 Brit. cwt. 2 qrs. 20 lbs. 13.75 oz.

The System Usuel, or Binary System, of the French, was introduced in 1812, for the accommodation of retail trade, to make a sort of compromise with the common people, who shewed an irreconcilable aversion to the innovations of the metrical or decimal system. It has the metrical standards for its basis; but their divisions, \&c. instead of being decimal, are chiefly binary, that is, they proceed principally by continually halving or doubling some of those standards, though partly according to other divisions of the old system; and instead of the metrical vocabulary, the names of ancient weights and measures are employed, annexing the term usuel to each. Thus, the toise usuelle = 2 metres = 78.74158 Brit. inches. The pied usuel is one sixth of the toise, and is subdivided into twelfths or inches, \&c. The aune usuelle = 12 decimetres = 1.2 metre = 47.2448 Brit. inches. The litron usuel = 1 litre = 1.7608 Brit. pint. The Boisseau usuel is one eighth of an hectolitre, and = 2.7512 Brit. gallons. The livre usuelle is half the kilogramme, and = 1.10243 lb. avoirdupois. The halves, quarters, eighths, \&c. of the most of the above are also in use. A summary of the ancient French system will be found in the pre-

Of the weights, 128 drachmes = 32 lot = 2 marks = 1 pound = 1.1145 lb. avoirdupois. The light pound is similarly divided, but only 1.0318 lb. avoirdupois, so that 108 of the light lbs. make only 100 of the heavy or centner weight. The Cologne mark, here reckoned = 3611 grains troy, is used for gold and silver. The Frankfort foot contains 11.42 Brit. inches, and the ell 21.54. The Brabant ell is generally used for Dutch goods, as the Paris aune for French. Eighty old or 90 new mass (each schoppen) are equal to 20 viertels = 1 ohm = 31.57 Brit. gall.: 16 seeters = 8 metzen = 4 simmers = 1 malter or achtel = 28 Brit. bushels. There are two sorts of pounds, the peso sottile lb., and the peso grosso lb. The latter is a tenth part heavier than the former, so that the cantaro of 100 lbs. peso sottile = 69.89 lbs. avoirdupois, and the cantaro of 100 lbs. peso grosso = 76.88. The latter is used for bulky commodities, and the former for gold, silver, and

ud for bulky commodities, and the former for gold, silver, and articles of small bulk. The palmo = 9-725 Brit. inches, and 203 palmi = 1 braccio. The canna is of three sorts : the canna piccola used by tradesmen = 9 palmi, the canna grossa of the merchants = 12 palmi, and the custom-house canna = 10 palmi. Of corn measure, 96 gombette = 8 quarti = 1 mina = 3-31 Brit. bushels. Also 100 pints = 2 barili = 1 mezzarola of wine = 32-67 imp. gall. ; and 64 quateroni = 4 quarti = 1 oil barile = 14-23 imp. gall.

Considerable diversity of weights and measures is found in the various states into which this extensive country is divided, and the most important of them will be found noticed separately. But not a few of their standards, though differing in amount, are similar in their multiples and subdivisions. The weights, 1024 hellers = 312 pfenninge = 128 quintins = 16 ounces = 2 marks = 1 commercial lb.; and 5760 troy = 288 scruples = 96 drams = 12 oz. = 1 apothecaries' lb.; 4352 eschen = 512 hellers = 256 pfennings = 64 quintins = 8 lots = 8 oz. = 1 Cologne mark = 3608 grains troy. By this gold and silver are weighed; and the fineness of gold is expressed in 24th parts or carats as in Britain, but the fineness of silver by 16th parts. For jewellery there is a carat of 3 171 grs. troy. Of measures, 144 inches = 12 feet = 6 ells = 2 klafters = 4 ruthe. The Rhineland foot used by surveyors = 12.36 Brit. inches; the long mile = 10,126 Brit. yds., the short mile = 6859, and the geographical mile = 801 Brit. yds. The Rhineland morgen of

Besides the British standards, the following Spanish also used: The arroba = 26 lbs. avoirdupois; the quintal = 100 = 101.75 lbs. avoirdupois; the corn fanega = 1.55 Brit. bushel; pipe of 117 gallons = 105 Brit. gall.; the liquid arroba = 7 Brit. gall.; and the wine gallon = 1.094 Brit. gall.

7 Brit. gall. ; and the wine gallon = 1-094 Brit. gall.
Goa.—The same as in Portugal, except that the candy of 20
sounds = 495 lbs. avoirdupois. This, in the measurement of grain, is
supposed to be nearly 14 Winchester bushels.

Greece.—The same as in France.
Guiana (British).—In general the

Greece.—The same as in France.
Guiana (British).—In general the same as in Britain; but the following, originally introduced by the Dutch, are also used: the lb. = 1-09 lbs. avoirdupois; and the ell of 26 inches = 27 British inches.

inches.
Grisma (Dutch).—In
ly upon the old system.

Hamburg.—Of the weights, 128 drachmes = 32 loths = 16 oz. = 2 marks = 1 lb. = 1-0682 lb. avoirdupois; 112 lbs. = 8 lispunds = 1 centner = 119.64 lbs. avoirdupois; and 5 centners = 2 shipfunds. Ten lbs. = 1 stone of wool or feathers = half a stone of flax. Of butter, 280 lbs. form the great, and 224 the small tonne. A pipe of oil is reckoned at 820 lbs. Gold and silver are weighed by the Cologne mark, and their fineness valued as explained under the head Germany. Of measures, 6 palms = 2 feet = 1 ell = 22.578 Brit. inches. The Rhineland foot of engineers and surveyors = 12.36 Brit. inches. The Brabant ell, commonly used for piece goods, = 27.585 Brit. inches. Also 160 quarters = 40 stubgens = 20 viertels = 5 elmers = 4 ankers = 1 ahm = 31.87 Brit. gall.; and 24 ankers or 6 ahms = 1 fuder. Of wine, 6 tierces = 4 oxhofts or hogsheads = 1 faas; but these are of various sizes. Of corn, 160 spints = 40 hintens = 20 faas = 10 scheffels = 1 wisp or wispel = 25 Brit. bushels; two wisps = 1 last of barley or oats = 7.25 Brit. quarters; and 3 wisps = 1 last of wheat or rye = 1 stock of barley or oats = 10.87 Brit. quarters.

Hanover.—The lb. = 1*073 lb. avoirdupois; 112 lbs. = 1 centner; and 20 lbs. = 1 stone of flax or 2 stones of wool. Also 600 lbs. = 51 stubgens = 2 tonnes of honey; and 3360 lbs. = 240 lispunds = 12 shipfunds = 1 last. The Cologne mark is used for gold, silver, and silk. Two feet = 1 ell = 22*91 Brit. inches; the mile = 6*5676 Brit. miles. The morgen land measure = 2*5625 Brit. roods. Also 80 kannen = 40 stubgens = 4 ankers = 2*5 eimers = 1 ahm = 34*24 Brit. gallons; and 6 ahms = 4 oxhofts = 1 fuder of wine. Of corn, 96 himtens = 16 malters = 2 wisps or wispels = 1 last = 82 Brit. bushels.

Hayti.—Port-au-Prince.—Principally the old system of France, together with the old English wine gallon.

Hesse-Cassel.—The lb. = 17*08 oz. avoirdupois. The ell = 22*59 Brit. inches; the acre = 5894 Brit. acre; the liquid ohm = 34*94 Brit. gall.; and 16 corn metzen = 4 himtens = 1 viertel = 4*42 Brit. bushels.

Brit. bushels.
Hesse-Darmstadt.—Two lbs. = 1 French kilogramme = 2-20486
lbs. avoirdu. Ten feet = 1 klafter = 2-5 metres = 8-2023 Brit. feet;
five ells = 3 metres = 3-2809 Brit. yds.; four morgens = 1 hec-
tare = 2-47114 Brit. acres. The liquid ohm = 160 litres = 35-22
Brit. gallons; the corn malter = 1-28 hectolitre = 3-522 Brit.
bushels.
Holland.—In 1820, the decimal system of France was introdu-
ced, but with the old Dutch nomenclature, as explained under the
head Belgium. Of the old weights still used in many places,
10-240 as = 320 engels = 16 oz. = 2 marks = 1 Dutch lb. troy

= 1'0851 lb. avoirdupois. There are also 16 oz. in the commercial lb. = 1'0893 lb. avoirdupois., and 100 commercial lbs. = 1 centner. The old Amsterdam foot = 11'15, and the ell = 27'08 Brit. inches; the Rhineland foot = 12'36, and the Flemish or Brabant ell = 27'58 Brit. inches. Nineteen Dutch leagues make one degree of the meridian.

meridian.
Japan.—The weights are almost the same as in China. The role of 52 Brit feet. Yet the measure of capacity have been compared.
inc = 6-25 Brit. feet; but the measures compared.

Java.—In this and the other Dutch colonies in India, the weights of China are ordinarily employed; but the pecul, instead of being exactly 133-33 lbs. avoirdupois, is = 135-625 lbs. Of rice, the coyang = 3581 lbs. avoirdupois, and the timbang = 10 sacks = 5 peculs = 678-125 lbs. The Dutch troy mark of 9 reals = 3798 Brit. troy grains, is used for gold and silver. The foot = 12-36, and the ell = 27-75 Brit. inches. The kanne, liquid measure, = 3282 Brit. gall., and 396 rands = 1 leaguer of arrack = 133-33 Brit. gall.; but 360 rands = 1 leaguer of wine.

India.—Grain is generally sold by weight, as also liquids, except wines and spirits, which at all the three presidencies are sold by British measures. Of the Bengal weights, 940 chittacks = 40 seers = 1 factory maund = 74.67 lbs. avoirdupois, or two thirds of a British cwt. The bazaar maund is one tenth greater than the factory, and similarly divided. Gold and silver are weighed by the new tola or sicca of 180 grains troy. Two cubits = 1 guz = 1 Brit. yd. The coss or mile of Bengal = 2000 yards, and 1600 sq. yds. = 20 cottahs = 1 bighah. At Bombay 40 seers = 1 maund = one quarter of a British cwt., and 20 maunds = 1 candy = 5 cwts., which in corn is counted at 24.5 Brit. bushels. At Madras, 320 pollams = 8 vis = 1 maund = 25 lbs. avoirdupois, and 20 maunds = 1 candy. The covid = 18.6 inches, but the British yard is used for cloth; and in land measure, 24 maunies = 1 cawney = 6400 Brit. sq. yds. Also 320 measures = 400 marcalis = 80 parahs = 1 garse


\section*{Source record}
\begin{description}
\item[Scan ID] \texttt{britannica:309f765b1ecc83e0a7f7762e44b4a927a7555443:train:104}
\item[Image file] \texttt{page-0104-9aaa833a4643.jpg}
\item[Image SHA-256] \texttt{9aaa833a46437aa0cc322c1d5f37b5a4d013906d6d40bfc13ca14d19046d6f3f}
\item[Dataset] NationalLibraryOfScotland/encyclopaedia\_britannica\_illustrated at revision \texttt{309f765b1ecc83e0a7f7762e44b4a927a7555443} (CC0-1.0)
\item[Extraction] model \texttt{nvidia/nemotron-3.5-parse}, raw reply \texttt{extract/raw/page-0104-9aaa833a4643.json}
\item[Raw reply SHA-256] \texttt{004b7e594ad00eb4a8b991c060453349d2039a3f9d4079cf5f4cf32109c46bef}
\item[Correction] model \texttt{nvidia/Nemotron-3.5-Super-VL}, block crops, raw replies \texttt{correct/raw/crop\_nothink/page-0104-9aaa833a4643-b*.json}; super\_vl: 18
\item[Running header] WEIGHTS AND MEASURES.
\end{description}
\end{document}
